% ===============================================================
% WEEKLY REPORT — copy this file to week02.tex, week03.tex, ...
% In Overleaf: use the file dropdown above the editor (or
% Menu > Settings > Main document) to pick which root to compile.
% ===============================================================
\documentclass[11pt,a4paper]{article}

\input{setup/project}
\input{setup/preamble}
\input{setup/weekly}

% No paragraph indentation; separate paragraphs with vertical space instead.
\setlength{\parindent}{0pt}
\setlength{\parskip}{0.6em}


\begin{document}

\weeklytitle{2}

\section{Week 2: 7--11 September 2026}


\sectionprogress

\textbf{Hemodynamics.} I read up on wall shear stress (WSS) in particular, and on fractional flow
reserve (FFR). It is well documented that regions of low or oscillatory WSS -- at bifurcations and
on the inner wall of bends -- are prone to plaque development \cite{chatzizisis2007ess,
asakura1990flow}. Far fewer studies relate the geometry of the coronary tree to its hemodynamics in
healthy arteries, and the role of flow in the onset of plaque remains underexplored for lack of
follow-up that starts from a healthy state \cite{zhang2024curvature}.

\textbf{Healthy population.} Most studies of coronary geometry are carried out in patients with
disease and ask where in the tree plaque appears, a question that requires plaque to be present. A
more interesting question for this thesis is whether the topology of the tree carries any signs of
risk \emph{before} disease develops. Phillip mentioned that the Copenhagen General Population Study
(CGPS; Herlev--{\O}sterbro) includes participants who were scanned twice, about ten years apart:
some had healthy trees at the first scan and had developed CAD by the second, while others remained
unchanged. Comparing these two groups could reveal geometric patterns that precede disease.

\textbf{Scope narrowed.} Rather than analyzing the whole coronary tree, I will begin the topology and
feature-extraction work with three vessels -- the left main (LM), left anterior descending (LAD) and
left circumflex (LCX) arteries -- for the reasons given in \Cref{sec:week2-major-vessel}.

\textbf{Features, focus on tortuosity.} Bifurcation angle is the one coronary feature with a
standardized definition, although the same angle can hide very different branch diameters.
Tortuosity is the opposite case. Its dominant measure, the tortuosity index $\tau = L/D$, cannot
tell apart vessels that bend in very different ways (\Cref{fig:week2-tortuosity-ambiguity}). This I will find a way to do better.  

\textbf{CAS-Net pipeline.} I went through the methods used in CAS-Net \cite{dong2023casnet} in more
depth, at a high level: data preparation and preprocessing, network architecture, training, inference
and evaluation metrics. As discussed, I will train CAS-Net from scratch. My contribution to improving
the segmentation for topological feature extraction will most likely focus on post-processing, since
the overlap with the reference annotations is already high, whereas the topology is not
\cite{bransby2026imagecasx}.

\textbf{Rigshospitalet.} I received an ID card for Rigshospitalet -- I am now an impostor medical
student.

\subsection{Questions}

\begin{itemize}
  \item We discussed that I should train CAS-Net from scratch. Would it be sufficient to train from
    scratch for only a few epochs (e.g.\ five) to confirm that the training pipeline runs, and then
    continue from the pretrained weights, rather than repeating the full training?
\end{itemize}

\newpage
\input{content/background/hemodynamics}

\section{Scope and State-of-the-Art Notes}
\input{content/Introduction/State of the art/segmentation}
\input{content/Introduction/State of the art/focus_on_major_vessel}
\label{sec:week2-major-vessel}
\input{content/Introduction/State of the art/features}
\input{content/Introduction/State of the art/healthy_cohort}

\newpage
\input{content/methods/pipeline/overview}
\input{content/methods/pipeline/data_preparation}
\input{content/methods/pipeline/model_architecture}
\input{content/methods/pipeline/training}
\input{content/methods/pipeline/inference}
\input{content/methods/pipeline/evaluation_metrics}

\newpage
\printbibliography

\end{document}
